\ifdefined\NoetherPaperTwentyTwoSFiveBodyOnly
\providecommand{\NoetherPaperTwentyTwoSFiveEnd}{}
\else
\documentclass[11pt]{article}
\usepackage[a4paper,margin=1.45cm]{geometry}
\usepackage[T1]{fontenc}
\usepackage[utf8]{inputenc}
\usepackage{lmodern}
\usepackage[french]{babel}
\usepackage{amsmath,amssymb,mathtools}
\usepackage[protrusion=true,expansion=false]{microtype}
\setlength{\parindent}{1.25em}
\setlength{\parskip}{0.10em}
\allowdisplaybreaks
\setlength{\emergencystretch}{12em}
\sloppy\hbadness=10000\vbadness=10000\hfuzz=2pt\vfuzz=10pt\raggedbottom
\makeatletter\renewcommand\footnotesize{\@setfontsize\footnotesize{7.5}{8.7}\selectfont}\makeatother
\providecommand{\Amod}{\mathfrak A}
\providecommand{\Bmod}{\mathfrak B}
\providecommand{\Gmod}{\mathfrak G}
\providecommand{\Mmod}{\mathfrak M}
\providecommand{\Nmod}{\mathfrak N}
\providecommand{\Rmod}{\mathfrak R}
\providecommand{\Smod}{\mathfrak S}
\providecommand{\mideal}{\mathfrak m}
\providecommand{\nideal}{\mathfrak n}
\providecommand{\gideal}{\mathfrak g}
\providecommand{\hideal}{\mathfrak h}
\providecommand{\jideal}{\mathfrak j}
\providecommand{\tideal}{\mathfrak t}
\providecommand{\aideal}{\mathfrak a}
\providecommand{\bideal}{\mathfrak b}
\providecommand{\bb}{\mathfrak b}
\providecommand{\cc}{\mathfrak c}
\providecommand{\zero}[1]{\equiv 0\pmod{#1}}
\providecommand{\Norm}{N}

\begin{document}
\providecommand{\NoetherPaperTwentyTwoSFiveEnd}{\end{document}}
\fi

% Unit: Paper 22 section 5. Direct French translation from the German final-audited source slice.
\setcounter{footnote}{13}

\section*{\S 5. Forme résultante et forme des diviseurs élémentaires d'un idéal de polynômes.}

Soient $x_{i+1}\ldots x_n$ adjoints à $P(u)$, de sorte que $\Gmod_{i-1}$ et $\Mmod_{i-1}$, et par conséquent aussi $\Gmod^*_{i-1}$ et $\Mmod^*_{i-1}$, passent en modules de formes linéaires où les grandeurs entières peuvent être regardées comme des polynômes d'\emph{une} indéterminée $x_i$. D'après le Théorème VII, les diviseurs élémentaires de $\Mmod_{i-1}$ qui sont différents de $1$, ainsi qu'au plus un nombre fini de diviseurs élémentaires $1$ de $\Mmod_{i-1}$, coïncident alors avec ceux de $\Mmod^*_{i-1}$. Le produit de ces diviseurs élémentaires, le $\rho$-ième diviseur déterminantiel de $\Mmod^*_{i-1}$, était égal au déterminant de la substitution de passage de $\Gmod^*_{i-1}$ à $\Mmod^*_{i-1}$, donc égal à $N(\Gmod^*_{i-1}\mid\Mmod^*_{i-1})$ d'après la Définition III. D'après (24), respectivement d'après la formule analogue qui résulte de (28) pour $i$ général, il apparaît que le produit de ces diviseurs élémentaires devient aussi égal au déterminant de la substitution de passage de $\Gmod_{i-1}$ à $\Mmod_{i-1}$; il doit donc être désigné par $N(\Gmod_{i-1}\mid\Mmod_{i-1})$. On obtient ainsi
\[
 N(\Gmod_{i-1}\mid\Mmod_{i-1})=N(\Gmod^*_{i-1}\mid\Mmod^*_{i-1})=R^{(i)}(x),
\]
où le polynôme $R^{(i)}(x)$, c'est-à-dire le plus grand commun diviseur, au sens polynomial, de tous les déterminants à $\rho$ lignes de $\Mmod^*_{i-1}$, peut être supposé entier et primitif en $x_{i+1}\ldots x_n$ et, par conséquent, comme diviseur de $C^{(i)}$, devient régulier en $x_i$. De même, le diviseur élémentaire le plus élevé $E^{(i)}(x)$ de $\Mmod_{i-1}$, respectivement de $\Mmod^*_{i-1}$, peut être supposé entier et primitif en $x_{i+1}\ldots x_n$, et par conséquent régulier en $x_i$. Cela conduit à la définition suivante.

\medskip
\noindent\textbf{Définition VI.} \emph{Le polynôme $R^{(i)}(x)$ est appelé la résultante de $i$-ième niveau de $\mideal$, et $E^{(i)}(x)$ le diviseur élémentaire de $i$-ième niveau. La résultante de $i$-ième niveau est donc la norme de l'idéal fondamental $\gideal_{i-1}$ par rapport à $\mideal$, conçue comme norme du module $\Gmod_{i-1}$ par rapport à $\Mmod_{i-1}$, où $x_{i+1}\ldots x_n$ sont adjoints à $P(u)$; de même, $E^{(i)}(x)$, sous la même adjonction, devient égal au quotient $\Mmod_{i-1}/\Gmod_{i-1}$. Comme forme résultante, respectivement forme des diviseurs élémentaires de $\mideal$, on désigne les produits suivants:}
\[
 R_\mideal=R^{(1)}R^{(2)}\cdots R^{(n)},\qquad
 E_\mideal=E^{(1)}E^{(2)}\cdots E^{(n)}.
\]

Pour la forme résultante et la forme des diviseurs élémentaires, on a:

\medskip
\noindent\textbf{Théorème VIII.} \emph{La forme résultante est divisible par la forme des diviseurs élémentaires; inversement, une puissance de $E_\mideal$ est divisible par $R_\mideal$. $E_\mideal$ et $R_\mideal$ sont divisibles par $\mideal$; plus généralement, $E^{(i)}\cdots E^{(n)}\gideal_{i-1}\equiv0(\mideal)$, et par conséquent $R^{(i)}\cdots R^{(n)}\gideal_{i-1}\equiv0(\mideal)$.}

En effet, comme la norme est divisible par le plus haut diviseur élémentaire et qu'inversement une puissance de ce plus haut diviseur élémentaire est divisible par la norme, on obtient cette divisibilité pour $R^{(i)}(x)$ et $E^{(i)}(x)$ après adjonction de $x_{i+1}\ldots x_n$. Mais puisque $R^{(i)}(x)$ et $E^{(i)}(x)$ sont supposés primitifs comme polynômes relativement à $x_{i+1}\ldots x_n$, la divisibilité vaut relativement à toutes les indéterminées $x_i,x_{i+1},\ldots,x_n$; ainsi vaut, d'après la définition de ces expressions relativement à toutes les indéterminées $x$, la divisibilité de $R_\mideal$ par $E_\mideal$, et d'une puissance de $E_\mideal$ par $R_\mideal$.

Le plus haut diviseur élémentaire $E^{(i)}(x)$ est en outre, d'après le Théorème VII et le Théorème II, défini, après adjonction de $x_{i+1}\ldots x_n$, comme le quotient $\Mmod_{i-1}/\Gmod_{i-1}$. Par conséquent, si l'on revient à des polynômes en toutes les indéterminées $x$, il existe pour tout
\begin{equation}
 G(x)\equiv0(\gideal_{i-1})
 \quad\hbox{un}\quad b^{(i+1)}\ne0,
 \quad\hbox{tel que}\quad b^{(i+1)}E^{(i)}(x)G(x)\equiv0(\mideal)
\tag{33}
\end{equation}
soit vérifié. Or $\gideal_0$ devient en particulier égal à l'idéal-unité, donc:
\[
 b^{(2)}E^{(1)}(x)\equiv0(\mideal),\quad b^{(2)}\ne0,
 \qquad\hbox{et donc}\qquad E^{(1)}(x)\equiv0(\gideal_1).
\]
En général, d'après (33), appliqué à
\[
 E^{(1)}(x)\cdots E^{(i-1)}(x)\equiv0(\gideal_{i-1}),
\]
il existe un $b^{(i+1)}\ne0$ tel que $b^{(i+1)}E^{(1)}\cdots E^{(i)}\equiv0(\mideal)$; donc $E^{(1)}\cdots E^{(i)}\equiv0(\gideal_i)$. Comme $\gideal_n=\mideal$, il vient ainsi
\begin{equation}
 E_\mideal=E^{(1)}\cdots E^{(n)}\equiv0(\mideal)
 \quad\hbox{et par conséquent}\quad
 R_\mideal=R^{(1)}\cdots R^{(n)}\equiv0(\mideal).
\tag{34}
\end{equation}

Si l'on fait parcourir à $G(x)$, dans (33), les polynômes en nombre fini d'une base d'idéal de $\gideal_{i-1}$, et si $c^{(i+1)}(x)$ désigne le produit des $b^{(i+1)}(x)$ qui leur correspondent, on obtient
\[
 c^{(i+1)}E^{(i)}(x)\gideal_{i-1}\equiv0(\mideal);
 \qquad c^{(i+1)}\ne0
 \quad\hbox{et donc}\quad E^{(i)}(x)\gideal_{i-1}\equiv0(\gideal_i),
\]
et de là, comme ci-dessus par répétition un nombre fini de fois:
\[
 E^{(n)}(x)\cdots E^{(i)}(x)\gideal_{i-1}\equiv0(\mideal)
\]
et par conséquent aussi
\[
 R^{(n)}(x)\cdots R^{(i)}(x)\gideal_{i-1}\equiv0(\mideal),
\]
ce qui démontre le Théorème VIII dans toutes ses parties.

Le Théorème VIII repose essentiellement sur des propriétés de la forme des diviseurs élémentaires; mais que la forme résultante ait une signification caractéristique pour $\mideal$ est montré par le théorème suivant.

\medskip
\noindent\textbf{Théorème IX.} \emph{Si $\nideal$ est un diviseur de $\mideal$ -- diviseur entendu au sens des idéaux -- et si les formes résultantes $R_\nideal$ et $R_\mideal$ coïncident, alors les idéaux $\nideal$ et $\mideal$ coïncident.}

Soient en effet $\gideal_0\ldots\gideal_{n-1},\gideal_n=\mideal$ et $\hideal_0\ldots\hideal_{n-1},\hideal_n=\nideal$ les idéaux fondamentaux de $\mideal$ et de $\nideal$. D'abord $\gideal_0$ et $\hideal_0$ deviennent égaux à l'idéal-unité, donc $\gideal_0=\hideal_0$. Supposons maintenant $\gideal_{i-1}=\hideal_{i-1}$, de sorte que $\Mmod_{i-1}$ et $\Nmod_{i-1}$ possèdent le même module fondamental $\Gmod_{i-1}$. L'hypothèse $R_\nideal=R_\mideal$ admet, après adjonction de $x_{i+1}\ldots x_n$, la forme
\[
 N(\Gmod_{i-1}\mid\Nmod_{i-1})=N(\Gmod_{i-1}\mid\Mmod_{i-1})
 \quad\hbox{ou encore}\quad
 N(\Gmod^*_{i-1}\mid\Nmod^*_{i-1})=N(\Gmod^*_{i-1}\mid\Mmod^*_{i-1});
\]
ici on a posé, conformément à (32): $\Nmod_{i-1}=(\Nmod^*_{i-1},\mathfrak C_{i-1})$, de sorte que $\Nmod^*_{i-1}$ devient donc aussi un module de formes linéaires en $\xi$, et que $\Gmod^*_{i-1}$ en est le module fondamental. En raison de la divisibilité de $\Mmod_{i-1}$ par $\Nmod_{i-1}$, qui entraîne la divisibilité de $\Mmod^*_{i-1}$ par $\Nmod^*_{i-1}$, il résulte alors du Théorème IV que, sous cette adjonction, les deux modules $\Mmod^*_{i-1}$ et $\Nmod^*_{i-1}$ coïncident, et par conséquent aussi $\Nmod_{i-1}$ et $\Mmod_{i-1}$. Mais l'adjonction de $x_{i+1}\ldots x_n$ signifie que l'on ajoute à $\Mmod_{i-1}$, respectivement à $\mideal$, tous les polynômes $G(x)$ tels que
\[
 b^{(i+1)}(x)G(x)\equiv0(\mideal);
 \qquad b^{(i+1)}\ne0
\]
soit vérifié, c'est-à-dire que, par l'adjonction, $\Mmod_{i-1}$, respectivement $\mideal$, passe en $\gideal_i$, et de façon correspondante $\nideal$ en $\hideal_i$; ainsi $\gideal_i=\hideal_i$, et par conséquent $\gideal_n=\hideal_n$, donc $\mideal=\nideal$ est démontré.

Enfin, le même principe de transfert, appliqué au Théorème V, donne encore:

\medskip
\noindent\textbf{Théorème X.} \emph{Soit $R^{(i)}=S^{(i)}T^{(i)}$ une décomposition de $R^{(i)}$ en deux polynômes $S^{(i)}$ et $T^{(i)}$ premiers entre eux -- au sens polynomial. Alors il existe un et un seul idéal $\jideal$, et de même un et un seul idéal $\tideal$, tous deux diviseurs de $\mideal$ avec même idéal fondamental de niveau $(i-1)$, $\gideal_{i-1}$, tels que $S^{(i)}$ devienne égal à $N(\Gmod_{i-1}\mid\Smod_{i-1})$ et $T^{(i)}$ égal à $N(\Gmod_{i-1}\mid\mathfrak T_{i-1})$. Les idéaux $\jideal$ et $\tideal$ sont définis comme la totalité des polynômes $S(x)$, respectivement $T(x)$, tels que}
\[
 s^{(i+1)}(x)T^{(i)}(x)S(x)\equiv0(\mideal),\quad s^{(i+1)}\ne0,
\]
\[
 t^{(i+1)}(x)S^{(i)}(x)T(x)\equiv0(\mideal),\quad t^{(i+1)}\ne0
\]
\emph{soit vérifié, et $\gideal_i$ devient le plus petit commun multiple de $\jideal$ et $\tideal$,}
\[
 \gideal_i=[\jideal,\tideal].
\]

Il reste seulement à remarquer que $s^{(i+1)},t^{(i+1)}$ peuvent être choisis fixes pour $\jideal$ et $\tideal$, comme produits des $s^{(i+1)},t^{(i+1)}$ correspondant aux éléments de base de $\jideal$, respectivement de $\tideal$, et que, comme on l'a remarqué ci-dessus, l'adjonction de $x_{i+1}\ldots x_n$ transforme $\mideal$ en $\gideal_i$. En particulier, la décomposition de $R^{(i)}$ peut être poursuivie jusqu'aux puissances de polynômes irréductibles -- facteurs primaires --, ce qui entraîne une représentation correspondante de $\gideal_i$ comme plus petit commun multiple.

\clearpage
\NoetherPaperTwentyTwoSFiveEnd
